AI data-loss-prevention (DLP) systems in Korean financial institutions must detect personal information whose categories are fixed by statute but whose surface forms in real text---call-centre transcripts, scanned forms, pasted tables---are anything but fixed. We introduce \textbf{Kiii-Kiii} (\emph{Korean Identifiers, Identifiability, and Ill-formed inputs}; dataset \kiii{}), a benchmark for Korean financial PII span extraction and de-identification whose name states its three axes. Its taxonomy has two tiers: a \emph{Legal} tier of 24 categories, each cited to an article of the Personal Information Protection Act, the Credit Information Act and its Enforcement Decree, the Real Name Financial Transactions Act, or the Electronic Financial Transactions Act, and an \emph{Identifiability} tier of 12 quasi-identifiers and behavioural attributes that enable profiling and hyper-personalised phishing. Documents are synthetic, long-context and multi-subject, and every identifier is generated by code with a valid checksum, then rendered through a controlled \emph{surface-form variation} axis (T0 canonical to T3 structural: dictated Korean numerals, cross-turn splitting, partial masking, agent read-back) calibrated to how commercial Korean speech-to-text engines actually emit numbers. We evaluate \todo{N} systems---regex/rule baselines, frontier APIs, open-weight models, and Korean-developed open LLMs---and find that \todo{headline finding: rules collapse at T2; Korean-native models help on X but not Y; partial-mask leakage}. The benchmark, taxonomy and generator are released \anonv{upon acceptance}{at \url{https://github.com/h-albert-lee/kiii-kiii-workspace}}.
